\section{From the trajectory record to the mass operator}
\label{vi:sec:operador-masa}

The construction distinguishes two operations. First, the dynamics produces
a trajectory record. Then a functional extracts integer
coordinates and evaluates them through the internal constants already obtained.
Preserving that order explains how mass can be a magnitude of
persistence without exhausting its structure.

\subsection{Chronological reading and oriented events}

Let \(d_1,\ldots,d_{108}\) be the digital reading of an enriched route.
In the stated chart, the action support is \(\{1,3,5,7\}\), and
\[
 n_A=\#\{t:d_t\in\{1,3,5,7\}\}.
\]
The \(12\) events of the \(9\)-step blocks are also preserved:
\[
 e_m=\sigma_m\#\{t:9m+1\le t\le9m+9,\ d_t\in\{2,4,6,8\}\},
 \quad 0\le m<12,
\]
where \(\sigma=(+,+,+,-,-,-,+,+,+,-,-,-)\). The support of this
extractor is an explicit part of the reading chart; its function must
not be confused with the ternary phase selector of TPK.

The distribution of events matters even when their sums
coincide. To see this, introduce
\[
 p_i=e_i+e_{i+6},\quad a_i=e_i-e_{i+6},\quad
 s_C=\sum_{i=0}^{5}p_i,\quad M_i=p_i-p_5\ (0\le i<5).
\]

\begin{proposition}[Integral reconstruction of events]
\label{vi:prop:eventos}
The transformation
\[
(e_0,\ldots,e_{11})\mapsto(s_C,M_0,\ldots,M_4,a_0,\ldots,a_5)
\]
is injective on \(\mathbb Z^{12}\). Its image is characterised by
\[
 s_C-\sum_{i=0}^4M_i\equiv0\pmod6,
 \qquad p_i\equiv a_i\pmod2\quad(0\le i<6),
\]
where \(p_5=(s_C-\sum M_i)/6\) and \(p_i=p_5+M_i\) for \(i<5\).
\end{proposition}
\begin{proof}
Adding the six expressions for \(p_i\) gives
\(s_C=6p_5+\sum_{i<5}M_i\). This recovers the \(p_i\), and then
\[
 e_i=\frac{p_i+a_i}{2},\qquad
 e_{i+6}=\frac{p_i-a_i}{2}.
\]
The congruences are necessary for these quantities to be integers and
sufficient for the formulas to define an integral preimage. Substitution
into the forward transformation recovers every coordinate.
\end{proof}

The decomposition \(12=1+5+6\) thus acquires an operational meaning:
the total, differences, and orientations recover the events.
Division by \(6\) belongs to this reconstruction. When the character
uses \(s_C/6\), \(s_C\) remains the integer sum; a coordinate already
divided is not divided a second time.

\subsection{Torsional memory and signature}

For every occurrence of \(9\), the reading preserves its predecessor.
Define, with cyclic indices,
\[
 j_2(d)=\begin{cases}0,&d=9,\\1,&d\in\{2,4,6,8\},\\-1,&d\in\{1,3,5,7\},\end{cases}
 \quad z_t=\mathbf1_{\{d_t=9\}}j_2(d_{t-1}).
\]
Then
\[
 T_z=\sum_{t=1}^{108}z_t,\quad
 C_z=\sum_{t\in\{27,54,81,108\}}z_t,\quad
 k_{\Delta_4}=T_z-2C_z.
\]
The first sum preserves accumulated memory, and the second distinguishes
corona closures. Removing the predecessors before evaluating these
expressions changes the object being read.

With \(\tau_m=\operatorname{sgn}(e_m)\), the remaining coordinates are
\[
 b_{120}=\sum_{a=0}^{3}\operatorname{sgn}(e_a+e_{a+4}+e_{a+8}),
 \qquad b_{\zeta}=\sum_{m=0}^{11}\tau_m\tau_{m+3},
\]
where the second index is modulo \(12\). The signature obtained is
\[
 L_{\mathrm{tick}}(\gamma)
 =\sigma(\gamma)=(n_A,s_C,k_{\Delta_4},b_{120},b_{\zeta})\in\mathbb Z^5.
\]
This map is deterministic in the stated chart. The signs and
correlations explain why composition must take place on the
record before extracting the \(5\) integers: adding signatures does not,
in general, represent the composition of two histories.

\subsection{Character and internal specialisation}

Let \(X=(X_A,X_C,X_4,X_{120},X_{270})\). The formal character is
\[
 \chi(\sigma;X)=\exp\!\left(n_AX_A+s_CX_C+k_{\Delta_4}X_4+
 b_{120}X_{120}+b_{\zeta}X_{270}\right).
\]
For each \(X\), it satisfies
\(\chi(\sigma+\sigma';X)=\chi(\sigma;X)\chi(\sigma';X)\)
because the exponent is linear in the signature. This is a law on signatures; it
does not assert that the extractor is additive on routes without preserving the interface.

Specialisation uses the internal coordinates constructed previously:
\[
 X_A=A_{\mathrm{rad}},\quad X_C=C^*_{\mathrm{rad}}/6,
 \quad X_4=\Delta_4,\quad X_{120}=s_{120},\quad
 X_{270}=\zeta_{270}:=135\Delta_4^2.
\]
The character is dimensionless and positive. Its effect is a homothety on
the mass line. The exponential form determines how the
coefficients act; their provenance is fixed by the preceding records and
constructions.

\subsection{Bilayer hierarchy and preservation of the two genealogies}

Write \(x=A_{\mathrm{rad}}\), \(y=C^*_{\mathrm{rad}}\). In the
chamber \(x>y>0\), the channels are
\[
 q_+=e^{-x-y},\qquad q_-=e^{-x+y},\qquad0<q_+<q_-<1.
\]
For an involution \(R\), with \(P_\pm=(I\pm R)/2\), the
two-channel operator is
\[
 T=q_+P_++q_-P_-.
\]
Its scalar and oriented moments are expressed as
\[
 c_n=q_-^n+q_+^n,\qquad s_n=q_-^n-q_+^n.
\]
Multiplying these expressions proves
\[
 c_n^2-s_n^2=4e^{-2nx},\quad
 c_{m+n}=\frac{c_mc_n+s_ms_n}{2},\quad
 s_{m+n}=\frac{s_mc_n+c_ms_n}{2}.
\]
Both sequences satisfy
\[
 X_{n+2}=c_1X_{n+1}-e^{-2x}X_n,\qquad X\in\{c,s\}.
\]
These relations show that the heights share the same operator.
A different collection of moments is not defined for every species.

\begin{proposition}[Semigroup of heights]
\[
 \langle90,120\rangle=\{90,120\}\cup\{30r:r\ge6\}.
\]
\end{proposition}
\begin{proof}
After division by \(30\), it suffices to study \(\langle3,4\rangle\).
The integers \(6,7,8\) are \(3+3,3+4,4+4\); adding \(3\) produces
every subsequent integer. Below \(6\), the only positive
elements are \(3,4\).
\end{proof}

The heights \(360=4\cdot90=3\cdot120\) illustrate the distinction between
equality of evaluation and memory of composition. The algebraic quotient
identifies the two exponents; the record preserves whether transport
used four stages of one type or three of the other.

The refined specialisation has exponent
\[
 \log\chi_{\mathrm{ref}}=n_Ax+\frac{s_C}{6}y+k_{\Delta_4}\Delta_4
 +b_{\zeta}\zeta_{270}+\sum_{h\in\langle90,120\rangle}N_hs_h,
 \qquad N_h=\eta_h+b_{120}\mathbf1_{\{h=120\}}.
\]
The pair \((\eta_{120},b_{120})\) is preserved even though evaluation uses
its sum. \(s_{270}\) and \(\zeta_{270}\) also remain
separate: the former is a bilayer moment, the latter a
quadratic expression of refinement.

\begin{proposition}[Convergence of the refined evaluation]
If \(\sum_h|N_hs_h|<\infty\), the refined character defines a finite
positive scalar, and its truncated evaluations converge to it.
\end{proposition}
\begin{proof}
The assumption gives absolute convergence of the exponent. Continuity of the
exponential permits passage to the limit; the exponential of a finite real number is
strictly positive.
\end{proof}

Convergence evaluates a sequence of coefficients already specified.
Selection of those coefficients by the route is a distinct operation
and must accompany every evaluation attributed to a species.

The notation \(\chi_{\mathrm{full}}\) retained in some sector
evaluation formulas denotes this same character
\(\chi_{\mathrm{ref}}\), with the same arguments and coefficients.
The change of symbol adds no further specialisation or tower.

\subsection{Displacement and autocorrelation readers}

The same route admits two readings of the return:
\[
 K_D=\left(D_{27}+D_{36},D_1,\frac{D_4-D_3}{2}\right),\qquad
 K_\Omega=(\Omega,54,P_6).
\]
The first map preserves displacement defects; the second, return
memory and autocorrelation. In the complete census of \(324\) orientations
there are \(14\) profiles of the first, \(9\) of the second, and \(40\) joint
pairs. The \(18\) coincidences have value \((80,54,6)\).
Equality of the two operators throughout the atlas is not inferred from this.

For the electronic route, the displacements are
\[
 (D_1,D_3,D_4,D_{27},D_{36})=(54,56,68,48,32).
\]
Substitution produces \(K_D=(80,54,6)\), coinciding with
\(K_\Omega\). Its exponent is
\[
 \kappa_e=80-54\alpha+6\Delta_4.
\]
The electronic trajectory is constructed in
\eqref{vi:dep:trayectoria-electron}; its displacement and
autocorrelation readers are evaluated, respectively, in
\eqref{vi:dep:KD} and \eqref{vi:dep:KOmega}. The coincidence
\(K_D=K_\Omega=(80,54,6)\) belongs to the beta-closure theorem in
Subsection~\ref{vi:dep:beta}. Here that central exception is preserved, and the
operation is extended to the atlas fibres.

\subsection{Operator on a fibre and dimensional section}

Let \(\mathscr H_M\) be the finite Hilbert space with orthonormal basis
\((e_\gamma)\) labelled by the routes of a multisection. For
\(r\in\{D,\Omega\}\), define the return operator by
\(B_M^re_\gamma=\kappa_\gamma^re_\gamma\), with real exponents;
it is therefore self-adjoint. Let \(M_M^{(0)}\) be the positive basal operator,
whose connection with the character is fixed below.
The dimensionless ratio of the
action sections is \(R_{\mathrm{act}}>0\). Define
\[
 M_M^{\beta,r}=R_{\mathrm{act}}^{B_M^r/2}
 M_M^{(0)}R_{\mathrm{act}}^{B_M^r/2}.
\]

\begin{proposition}[Positivity and scalar restriction]
The operator \(M_M^{\beta,r}\) is positive. If both operators are
diagonal in the route basis, each basal value \(m_\gamma^{(0)}\)
is transformed into
\[
 m_\gamma^\beta=m_\gamma^{(0)}R_{\mathrm{act}}^{\kappa_\gamma}.
\]
\end{proposition}
\begin{proof}
Let \(S=R_{\mathrm{act}}^{B_M^r/2}\), self-adjoint and invertible.
For every \(v\),
\(\langle v,SM_M^{(0)}Sv\rangle=\langle Sv,M_M^{(0)}Sv\rangle\ge0\).
In the diagonal basis, the two factors \(S\) multiply each value
by \(R_{\mathrm{act}}^{\kappa_\gamma/2}\); their product gives the formula.
\end{proof}

The absolute mass section is obtained from action and the clock:
\[
 \omega_0=\frac{2\pi}{108t_0},\qquad
m_{\mathrm{clk}}=\hbar_{\mathrm{ret}}\omega_0c_{\mathrm{int}}^{-2}.
\]
In the chart normalised relative to the electron, the preceding basal operator is
\[
 M_M^{(0)}e_\gamma
 =m_e^{(0)}\,
   \chi_{\mathrm{ref}}(\sigma_\gamma-\sigma_e,
                       \eta_\gamma-\eta_e)e_\gamma.
\]
Here \(m_e^{(0)}\) is the realisation of the basal electronic
coordinate, with its prefactor, refinement, and dimensional basis.
If \(E_e^{(0)}=\varepsilon_e^{(0)}\mathcal U_E\), then
\[
 m_e^{(0)}=\varepsilon_e^{(0)}\mathcal U_Ec_{\mathrm{int}}^{-2}
 =b_e\,m_{\mathrm{clk}},\qquad
 b_e=\varepsilon_e^{(0)}\frac{\mathcal U_E}{E_{\mathrm{clk}}}.
\]
In the chart \(\mathcal U_E=E_{\mathrm{clk}}\), we obtain
\(b_e=\varepsilon_e^{(0)}\); in another chart, the stated ratio
of bases is retained. If the tower coordinates are already
relative to the electron, a second subtraction of \(\eta_e\) is omitted.
This equality links the operator to the signature of each route and explicitly
preserves electronic normalisation. A different choice of
reference section must also transport that normalisation.
The scale of the reduced Planck constant belongs to this dimensional
construction. The character and return factor subsequently act as
scalars. The electron preserves its structure: it is not identified with the
chronological quantum \(m_{\mathrm{clk}}\).

For two routes in the same absolute chart,
\[
 \frac{m_Y^\beta}{m_X^\beta}
 =\chi_{\mathrm{ref}}(\sigma_Y-\sigma_X,\eta_Y-\eta_X)
 R_{\mathrm{act}}^{\kappa_Y-\kappa_X}.
\]
The common dimensional section cancels. The difference of the return
exponents does not cancel unless it is zero. This identity preserves the
distinction between the absolute scale and the relative structure of a collection.

\subsection{The null propagation sector}

The exponential takes strictly positive values. The null chart is
constructed before that character, in the split realisation
\[
 \mathcal A_{\mathrm{sp}}=\mathbb R[R]/(R^2-1),\quad
 z=r_+P_++r_-P_-,\quad N(z)=r_+r_-.
\]
The directions \(\mathbb RP_+\) and \(\mathbb RP_-\) have zero
norm. Propagation transported in one of them belongs to that sector;
the bilayer material reading uses the coupling of both. In this way,
the photon and the gluonic sector retain their place in the theory even though
they do not form a row of the positive character. The realisation of their
representations is presented alongside the corresponding sectors.
